\documentclass[10pt,a4paper]{amsart}
\usepackage[margin=19mm]{geometry}
\usepackage{amsmath,amssymb,mathtools}
\usepackage[scr=rsfs,cal=euler]{mathalfa}
\usepackage{xfrac}
\usepackage[unicode,hidelinks,bookmarks=false]{hyperref}
\newcommand{\ie}{i.e.\ }
\newcommand{\cf}{cf.\ }
\newcommand{\resp}{respectively\ }
\newcommand{\Subsection}{\subsection}
\providecommand{\nobreakdash}{\nobreak\hbox{-}}
\setlength{\emergencystretch}{2em}
\multlinegap0pt
\usepackage{latexsym,cancel,rotating}
\usepackage{mathrsfs,color,fancyhdr}
\usepackage{xy}
\usepackage{verbatim}
\usepackage{xy}
\usepackage{rotating}
\usepackage{tikz-cd}
\usepackage{setspace}
\newcommand{\margin}[1]{\marginpar{\tiny #1}}
\newtheorem{theorem}{Theorem}[section]
\newtheorem{lemma}[theorem]{Lemma}
\newtheorem{corollary}[theorem]{Corollary}
\newtheorem{definition}[theorem]{Definition}
\newtheorem{proposition}[theorem]{Proposition}
\newtheorem{conjecture}[theorem]{Conjecture}
\newtheorem{remark}[theorem]{Remark}
\newtheorem{problem}[theorem]{Problem}
\newtheorem{example}[theorem]{Example}
\newtheorem{question}[theorem]{Question}
\newtheorem{thm}{Theorem}
\newcommand{\dleb}{\left[\!\!\left[}
\newcommand{\leb}{\left[}
\newcommand{\drib}{\right]\!\!\right]}
\newcommand{\rib}{\right]}
\def\lan{\langle}    \def\ran{\rangle}
\def\lr#1{\langle #1\rangle}
\def\lrs#1{< #1>}
\def\az{\alpha}  \def\thz{\theta}
\def\bz{\beta}  \def\dt{\Delta}
\def\gz{\gamma}
\def\ggz{\Gamma}  \def\ooz{\Omega}\def\ssz{\Sigma}
\def\sz{\sigma}
\def\oz{\omega}
\def\vez{\varepsilon}  \def\ez{\epsilon}
 \def\dz{\delta}
\def\pz{\partial}
\def\llz{\Lambda}
\def\lz{\lambda}
\def\cM{{\mathcal M}}\def\cP{{\mathcal P}} \def\fkR{{\frak R}}
\def\cC{{\mathcal C}} \def\cD{{\mathcal D}} \def\cX{{\mathcal X}} \def\cY{{\mathcal Y}}  \def\cL{{\mathcal L}}
\def\cH{{\mathcal H}}   \def\cB{{\mathcal B}} \def\cS{{\mathcal S}} \def\fkS{{\frak S}}
\def\cZ{{\mathcal Z}} \def\cO{{\mathcal O}} \def\cE{{\mathcal E}}
\def\cI{{\mathcal I}}  \def\cJ{{\mathcal J}}  \def\cQ{{\mathcal Q}}  \def\cR{{\mathcal R}}
\def\cA{{\mathcal A}}  \def\scrP{{\scr P}}  \def\scrL{{\scr L}}
\def\cK{{\mathcal K}}
\def\bbN{{\mathbb N}}  \def\bbZ{{\mathbb Z}}  \def\bbQ{{\mathbb Q}}
\def\bbR{{\mathbb R}}  \def\bbT{{\mathbb T}}  \def\bbF{{\mathbb F}}
\def\bbC{{\mathbb C}}  \def\bbE{{\mathbb E}}  \def\bbP{{\mathbb P}}
\def\bbK{{\mathbb K}}
\def\fh{\mathfrak h}     \def\tpi{{\tilde \pi}}    \def\tlz{{\tilde \lz}}
\def\iz{\iota}   \def\tiota{{\tilde \iota}}   \def\vphi{\varphi}
\def\dpeq{{\prec\!\!\!\cdot}}        \def\bfd{{\bf d}} \def\bfU{{\bf U}}
\def\bC{{\bf C}}  \def\barP{{\bar P}} \def\barH{{\bar H}}
\def\peq{\preccurlyeq}  \def\leq{\leqslant}  \def\geq{\geqslant}
\def\ged{\geqslant} \def\led{\leqslant}
\def\lra{\longrightarrow}   \def\skrel{\stackrel}
\def\lmto{\longmapsto}   \def\overQ{{\overline Q}}
\def\ra{\rightarrow}
\def\Hom{\mbox{\rm Hom}}  \def\ap{\mbox{\rm ap}}
\def\op{\mbox{\rm op}}       \def\im{\mbox{\rm Im}\,}
\def\Ext{\mbox{\rm Ext}\,}   \def\Ker{\mbox{\rm Ker}\,}
\def\dim{\mbox{\rm dim}\,}   \def\End{\mbox{\rm End}}
\def\Aut{\mbox{\rm Aut}}
\def\udim{{\mathbf dim\,}} \def\oline{\overline}
\def\mod{\mbox{\rm mod}\,} \def\mx{\mbox{\rm max}} \def\tr{\mbox{\rm tr}}
\def\rad{\mbox{\rm rad}\,}  \def\top{\mbox{\rm top}\,}
\def\prep{\mbox{\rm Prep}\,}  \def\prei{\mbox{\rm Prei}\,}
\def\reg{\mbox{\rm Reg}\,}
\def\fkm{{\frak m}}  \def\fku{{\frak u}}
\def\fkb{{\frak b}}  \def\fkc{{\frak c}}
\def\fkL{{\frak L}}  \def\fkp{{\frak p}}
\def\fkO{{\frak O}}
\def\proclaim#1{\vskip.3cm\noindent{\bf#1.}\quad\it}
\def\endproclaim{\par\vskip.3cm\rm}
\def\bfb{{\bf b}}
\def\bfV{{\mathbf V}}
\def\bfk{{\mathbf k}}
\def\bfH{{\mathbf H}} \def\bfP{{\mathbf P}} \def\bfa{{\bf a}}
\def\bfc{{\boldsymbol c}}  \def\bfd{{\bf d}}
\def\tu{{{\widetilde u}}}
\def\dg{{\rm deg}}
\def\blz{{|\lz\rangle}}
\def\bmu{{|\mu\rangle}}
\def\bnu{{|\nu\rangle}}
\def\wtPi{{\widetilde\Pi}}
\def\bempty{{|\emptyset\rangle}}
\def\supp{{\rm supp}}
\def\diam{{\diamond}}
\def\bfp{{\bf p}}
\def\scrF{{\scr F}}
\def\oln{\overline}
\def\fkM{{\frak M}}
\def\bfcH{\boldsymbol{\mathcal H}}
\def\bfcC{\boldsymbol{\mathcal C}}
\def\scrD{\boldsymbol{\mathcal D}}
\def\bfm{{\bf m}}
\def\wt{\widetilde}
\def\bdc{\boldsymbol{ c}}
\def\bdx{\boldsymbol{ x}}
\def\bdz{\boldsymbol{ z}}
\def\Rep{{\rm Rep}}
\def\wh{\widehat}
\def\dHallp{{\scrD(\dt)^{\geq 0}}}
\def\dHalln{{\scrD(\dt)^{\leq 0}}}
\def\bfd{{\bf d}}
\def\fka{{\mathfrak a}}
\def\ot{\otimes}
\def\afH{{\widehat {\bf H}}}
\def\bfi{{\bf i}}
\def\bfn{{\bf n}}
\newcommand \ud{{\underline{d}}}
\newcommand \bd{{\bf d}}
\newcommand\bbn{{\mathbb N}}
\newcommand\bbz{{\mathbb Z}}
\newcommand\bbq{{\mathbb Q}}
\newcommand\bbc{{\mathbb C}}
\newcommand\bba{{\mathbb A}}
\newcommand\fg{\mathfrak g}
\newcommand\fn{\mathfrak n}
\newcommand\tg{{\tilde \fg}}
\newcommand\tmu{{\tilde \mu}}
\newcommand\tA{{\tilde A}}
\newcommand\taz{{\tilde \alpha}}
\newcommand\trho{{\tilde \rho}}
\newcommand\te{{\tilde E}}
\newcommand\bi{{\bf 1}}
\newcommand\tnu{{\tilde \nu}}
\newcommand\mood{{\text{mod}}}
\newcommand\mi{{\mathcal{I}}}
\newcommand\mo{{\mathcal{O}}}
\newcommand\mq{{\mathcal{Q}}}
\newcommand\mr{{\mathcal{R}}}
\newcommand\ml{{\mathcal{L}}}
\newcommand\mh{{\mathcal{H}}}
\newcommand\udd{{\underline{\bf d}}}
\newcommand\ue{{\underline{e}}}
\newcommand\uee{{\underline{\bf e}}}
\newcommand\gt{{\langle G_{\bd},T_{\bd}\rangle}}
\begin{document}
\title[Notes on Chevalley Groups and Root Category IV: Iwasawa Deco]{Notes on Chevalley Groups and Root Category IV: Iwasawa Decomposition and Total Positivity}
\author{Buyan Li}
\date{}
\maketitle
\noindent\emph{Source and licence.} Notes on Chevalley Groups and Root Category IV: Iwasawa Decomposition and Total Positivity. Version arXiv:2607.23697v1 (CC BY 4.0). This material is distributed under the Creative Commons Attribution 4.0 International licence, http://creativecommons.org/licenses/by/4.0/, as recorded in the arXiv OAI licence metadata for this exact version and shown on the abstract page https://arxiv.org/abs/2607.23697v1. Full rights notice: licenses/arxiv-2607.23697-CC-BY-4.0.txt.\par\smallskip
\noindent\textit{Editorial scope.} Counted unit: the original \texttt{theorem} environment \texttt{-} at source lines 585--588 together with its complete proof at lines 590--607; the delivered document keeps source lines 307--608 so that every referenced label and every citation resolves inside it. Native editable source is the paper's own concatenated TeX; the AMS wrapper is editorial formatting only. No publisher class, upstream script or unrelated artwork is redistributed.
\section{Cartan Decomposition and Iwasawa Decomposition}
In this section, we obtain the Cartan decomposition and the Iwasawa decomposition for the simple Lie algebra and the Chevalley group from the root category $\mathcal{R}$, with the field $\mathbb{K}$ taken to be $\mathbb{C}$.
For the classical results of these decompositions, see reference \cite{Knapp}; however, note that in our case, these decompositions are more intuitional and simpler.

From now on, we denote $\mathfrak{g}=\mathfrak{g}_{\mathbb{C}}$ and $\mathbf{G}=\mathbf{G}(\mathcal{R},\mathbb{C})$.

\subsection{Cartan Decomposition}

In \cite{notes2}, we have essentially obtained the Cartan decomposition for both the Lie algebra and the Chevalley group, though not stated explicitly.
We now present it in full, supplementing the necessary details.

Recall that the shift functor $T$ of the root category $\mathcal{R}$ induces an $\mathbb{R}$-linear involution $\theta$ of $\mathfrak{g}$, given by 
\begin{align*}
    \theta(H_X)=H_{TX}, \quad \theta(u_X)=u_{TX},
\end{align*}
for each $X\in \operatorname{ind}\mathcal{R}$, and $\theta$ maps the imaginary unit $i$ to $-i$. 
We call $\theta$ the Cartan involution of $\mathfrak{g}$.

Let 
\begin{align*}
    \mathfrak{r}=\{ X\in \mathfrak{g}| \theta(X)=X \}
\end{align*}
be the $\theta$-fixed points of $\mathfrak{g}$, which is a real Lie algebra.
By an approach analogous to that in \cite{1997ROOT}, we construct a $\mathbb{Z}$-form of the real Lie algebra $\mathfrak{r}$ from the root category $\mathcal{R}$.
As a result, we obtain a $\mathbb{Z}$-basis of $\mathfrak{r}$. 
Precisely, if we choose a subset $\Delta$ of $\operatorname{ind}\mathcal{R}$ such that $\{H_X,X\in \Delta\}$ form a basis of $\mathcal{K}$, (for example, the set of simple objects in a complete hereditary subcategory), then 
\begin{align*}
    \{i\frac{H_X}{d(X)},X\in \Delta\}\bigcup \{u_X+u_{TX},i(u_X-u_{TX}),X\in \operatorname{ind}\mathcal{R}\}
\end{align*}
form a $\mathbb{Z}$-basis of $\mathfrak{r}$. 
Here $d(X)=\operatorname{dim}_k\operatorname{End}_{\mathcal{R}}X$ for any $X\in \operatorname{ind}\mathcal{R}$.
Moreover, we have the following result.

\begin{proposition}
    \cite[Proposition 30]{notes2}
    The real Lie algebra $\mathfrak{r}$ is a compact real form of $\mathfrak{g}$.
\end{proposition}

Since $\theta$ is an involution, it has eigenvalues $\pm 1$.
Let $\mathfrak{p}$ be the eigenspace corresponding to eigenvalue -1. 
Then $\mathfrak{g}$ has an eigenspace decomposition
\begin{align*}
    \mathfrak{g}=\mathfrak{r}\oplus \mathfrak{p}.
\end{align*}
This is called the Cartan decomposition of $\mathfrak{g}$.
Note that $\mathfrak{r}$ is $\mathbb{R}$-spanned by $i\frac{H_X}{d(X)}$, $u_X+u_{TX},i(u_X-u_{TX})$, $X\in \operatorname{ind}\mathcal{R}$, and $\mathfrak{p}$ is $\mathbb{R}$-spanned by $\frac{H_X}{d(X)},u_X-u_{TX},i(u_X+u_{TX})$, $X\in \operatorname{ind}\mathcal{R}$.

Now we consider the Cartan decomposition of the Chevalley group $\mathbf{G}$.
Recall that in \cite{notes2} we let $\mathbf{R}$ be the subgroup of $\operatorname{Aut}(\mathfrak{r})$ generated by $\operatorname{exp}(t\operatorname{ad}i\frac{H_X}{d(X)})$, $\operatorname{exp}(t\operatorname{ad}(u_X+u_{TX})),\operatorname{exp}(t\operatorname{ad}i(u_X-u_{TX}))$, with $X\in \operatorname{ind}\mathcal{R}$ and $t\in \mathbb{R}$, and let $\mathbf{K}$ be the identity component of $\mathbf{R}$.
As a corollary of \cite[Remark 37]{notes2}, we actually have $\mathbf{K}=\mathbf{R}$.

\begin{proposition}
    \cite[Proposition 36]{notes2}
    The group $\mathbf{K}$ is a maximal compact subgroup of $\mathbf{G}$. 
    In fact, it coincides with the inner automorphism group  $\operatorname{Int}(\mathfrak{r})$.
\end{proposition}

We define an involution $\Theta$ on $\mathbf{G}$.
Let $\Theta:\mathbf{G}\rightarrow \mathbf{G}$ be a group automorphism, mapping $E_{[X]}(t)$ to $E_{[TX]}(\overline{t})$, for any $X\in \operatorname{ind}\mathcal{R}$ and $t\in \mathbb{C}$.
Here $\overline{t}$ means taking the conjugate of $t$.
It can be easily shown that $\Theta^2=1$, $\Theta(n_{[X]})=n_{[TX]}$, and $\Theta(h_{[X]}(t))=h_{[TX]}(\overline{t})$, for any $X\in \operatorname{ind}\mathcal{R}$ and $t\in \mathbb{C}$.
Moreover, for any $t\in \mathbb{C}$, since
\begin{align*}
     \Theta(\operatorname{exp}(t\operatorname{ad}u_X))&=\Theta(E_{[X]}(t))  =E_{[TX]}(\overline{t})=\operatorname{exp}(\overline{t}\operatorname{ad}u_{TX})=\operatorname{exp}(\operatorname{ad}\theta(tu_X)),\\
     \Theta(\operatorname{exp}(t\operatorname{ad}\frac{H_X}{d(X)}))&=\Theta(h_{[X]}(e^{-t}))=h_{[TX]}(e^{-\overline{t}})=\operatorname{exp}(\overline{t}\operatorname{ad}\frac{H_{TX}}{d(X)})=\operatorname{exp}(\operatorname{ad}\theta(t\frac{H_X}{d(X)})),
\end{align*}
we have $d\Theta=\theta$.

Denote the set of $\Theta$-fixed points in $\mathbf{G}$ by $\mathbf{G}^{\Theta}$.
Then the Lie algebra of $\mathbf{G}^{\Theta}$ is $\mathfrak{r}$.

\begin{lemma}
    \cite[Lemma 31]{notes2}
    For any $t\in \mathbb{R}, X\in \operatorname{ind}\mathcal{R}$, we have 
    \begin{align*}
       & \operatorname{exp}(t\operatorname{ad}(u_X+u_{TX}))=E_{[X]}(\operatorname{tan}t)h_{[X]}((\operatorname{cos}t)^{-1})E_{[TX]}(\operatorname{tan}t),\\
        & \operatorname{exp}(t\operatorname{ad}i(u_X-u_{TX}))=E_{[X]}(i\operatorname{tan}t)h_{[X]}((\operatorname{cos}t)^{-1})E_{[TX]}(-i\operatorname{tan}t).
    \end{align*}
\end{lemma}

Using the definition and the above lemma, we notice that $\mathbf{K}$ is contained in $\mathbf{G}^{\Theta}$.
We denote the identity component by $(-)^{\circ}$, and denote the Lie algebra of the Lie group by $\operatorname{Lie}(-)$.

\begin{proposition}
    The group $\mathbf{K}$ coincides with $\mathbf{G}^{\Theta}$.
\end{proposition}

\begin{proof}
    The proposition follows easily from the facts that $\mathbf{K}=\operatorname{Int}(\mathfrak{r})=\operatorname{Aut}(\mathfrak{r})^{\circ}$, $\operatorname{Lie}(\mathbf{G}^{\Theta})=\mathfrak{r}=\operatorname{Lie}(\mathbf{K})$, and $\mathbf{G}^{\Theta}$ is a connected, closed subgroup of $\operatorname{Aut}(\mathfrak{r})$.
\end{proof}

Recall the definition of an invariant symmetric bilinear form on $\mathfrak{g}$ arising from $\mathcal{R}$.

\begin{definition}\cite[Definition 29]{notes2}\label{(,)}
Let $(,):\mathfrak{g} \times \mathfrak{g} \rightarrow \mathbb{C}$ be a non-degenerate symmetric bilinear form on $\mathfrak{g}$, satisfying the following properties:

(1) For any $x,y,z\in \mathfrak{g}$, we have 
\begin{align*}
    ([x,y],z)=(x,[y,z]).
\end{align*}
    
(2) For any $X,Y\in \mathcal{R}$, we have 
\begin{align*}
    (H_X,H_Y)=\sum_{Z\in \operatorname{ind}\mathcal{R}}(H_X|H_Z)(H_Y|H_Z),
\end{align*}
where $(-|-)$ is the symmetric Euler form on $\mathcal{K}$.

(3) For any $X,Y\in \operatorname{ind}\mathcal{R}$, $(H_X,u_Y)=0$, i.e. $(\mathcal{K},\mathcal{N})=0$.

(4) For any $X,Y\in \operatorname{ind}\mathcal{R}$, if $X\ncong TY$, then $(u_X,u_Y)=0$. 
If $X\cong TY$, then 
\begin{align*}
    (u_X,u_{TX})= -4+\sum_{Y\in\operatorname{ind}\mathcal{R}}(\sum_{L\in\operatorname{ind}\mathcal{R}} \gamma_{TX,Y}^L\gamma_{X,L}^Y),
\end{align*}
where integers $\gamma_{MN}^P$ are evaluations of the Hall polynomials and are used to define the Lie bracket of $\mathfrak{g}$ by Peng and Xiao. 
\end{definition}

\begin{lemma}
    For any $x\in \mathbf{G}$ and $Y,Z\in \mathfrak{g}$, we have 
    \begin{align*}
        (xY,xZ)=(Y,Z).
    \end{align*}
\end{lemma}

\begin{proof}
    It suffices to prove for generators $E_{[X]}(t)$ of the Chevalley group.
    It follows from Definition \ref{(,)}(1) that $((\operatorname{ad}y)(x),z)=-(x,(\operatorname{ad}y)(z))$, for any $x,y,z\in \mathfrak{g}$.
    Now for $X\in \operatorname{ind}\mathcal{R},t\in \mathbb{C},Y,Z\in \mathfrak{g}$, we have 
    \begin{align*}
        (E_{[X]}(t)Y,Z)&=((1+t\operatorname{ad}u_X+\frac{t^2}{2!}(\operatorname{ad}u_X)^2+\cdots)Y,Z)\\
       & =(Y,(1-t\operatorname{ad}u_X+\frac{(-t)^2}{2!}(\operatorname{ad}u_X)^2+\cdots)Z)=(Y,E_{[X]}(-t)Z),
    \end{align*}
    and the lemma is proved.
\end{proof}

By definition, it is easy to check that $(\theta (X),\theta (Y))=(X,Y)$ for any $X,Y\in \mathfrak{g}$.

Now we define a bilinear form $(-,-)_{\theta}$ on $\mathfrak{g}$ by 
\begin{align*}
    (X,Y)_{\theta}=-(X,\theta (Y)),
\end{align*}
for any $X,Y\in \mathfrak{g}$.
Then 
\begin{align*}
    (Y,X)_{\theta}=-(Y,\theta(X))=-(\theta(Y),X)=-(X,\theta (Y))=(X,Y)_{\theta}
\end{align*}
and $(-,-)_{\theta}$ is symmetric.

\begin{lemma}
    The symmetric bilinear form $(-,-)_{\theta}$ is positive definite.
\end{lemma}

\begin{proof}
    By definition, we notice that $(-,-)_{\theta}$ coincides with $(-,-)$ when restrict to $\mathcal{K}$, and $(\mathcal{K},\mathcal{N})_{\theta}=0$.
    For any $X,Y\in \operatorname{ind}\mathcal{R}$, if $X\ncong Y$, then 
    \begin{align*}
        (u_X,u_Y)_{\theta}=-(u_X,u_{TY})=0.
    \end{align*}
    Otherwise,
    \begin{align*}
        (u_X,u_X)_{\theta}=-(u_X,u_{TX})>0,
    \end{align*}
    which is shown in the proof of \cite[Proposition 30]{notes2} by properties of $\gamma$.
    Thus $(-,-)_{\theta}$ is positive definite.
\end{proof}

We regard $\mathfrak{g}$ as a real Lie algebra, and then $(-,-)_{\theta}$ is an inner product on it.

\begin{lemma}
    The Chevalley group $\mathbf{G}$ coincides with $\operatorname{Int}(\mathfrak{g})$.
\end{lemma}

\begin{proof}
    Since $\operatorname{Int}(\mathfrak{g})$ is generated by $\operatorname{exp}(\operatorname{ad}X), X\in \mathfrak{g}$, $\mathbf{G}$ is contained in $\operatorname{Int}(\mathfrak{g})$.
    On the other hand, the Lie algebras $\operatorname{Lie}(\mathbf{G})=\mathfrak{g}=\operatorname{Lie}(\operatorname{Int}(\mathfrak{g}))$, so $\mathbf{G}^{\circ}=\operatorname{Int}(\mathfrak{g})^{\circ}=\operatorname{Int}(\mathfrak{g})$, and the lemma follows.
\end{proof}

\begin{theorem}
    The map $\mathbf{K}\times \mathfrak{p}\rightarrow \mathbf{G}, (k,X)\mapsto k\operatorname{exp}(\operatorname{ad}X)$ is a diffeomorphism.
\end{theorem}

\begin{proof}
    We first show that this map is surjective.
    For any $x\in \mathbf{G}$, we consider the element $y=\Theta(x)^{-1}x$.
    Then $\Theta(y)=y^{-1}$.
    Moreover, for any $X,Y,Z\in \mathfrak{g}$, we have 
    \begin{align*}
        (yX,Y)_{\theta}&=-(yX,\theta(Y))=-(X,y^{-1}\theta(Y))\\
        &=-(X,\Theta(y)\theta(Y))=-(X,\theta(yY))=(X,yY)_{\theta},
    \end{align*}
and
\begin{align*}
    (yZ,Z)_{\theta}=-(\Theta(x)^{-1}xZ,\theta(Z))=-(xZ,\theta(xZ))=(xZ,xZ)_{\theta}\geq 0.
\end{align*}
Here $ (yZ,Z)_{\theta}=0$ if and only if $xZ=0$, if and only if $Z=0$.
As a result, the element $y$, viewed as an automorphism of $\mathfrak{g}$, is diagonalizable and all its eigenvalues are positive real numbers.
We choose a basis of $\mathfrak{g}$ made of eigenvectors of $y$ and denote the eigenvalues of $y$ by $\lambda_1,\cdots,\lambda_m$.
With respect to this basis, $y$ becomes a diagonal matrix $\operatorname{diag}\{\lambda_1,\cdots,\lambda_m\}$.
For any $t\in \mathbb{R}$, define $y^t=\operatorname{diag}\{\lambda_1^t,\cdots,\lambda_m^t\}$, which is an automorphism of $\mathfrak{g}$.
The elements $\{y^t,t\in \mathbb{R}\}$ form a one-parameter subgroup of $\operatorname{Aut}(\mathfrak{g})$, and thus contained in $\operatorname{Aut}(\mathfrak{g})^{\circ}=\operatorname{Int}(\mathfrak{g})=\mathbf{G}$.
So $\{y^t,t\in \mathbb{R}\}$ is a one-parameter subgroup of $\mathbf{G}$, and there exists some $X\in \mathfrak{g}$ such that $y^t=\operatorname{exp}(t\operatorname{ad}X)$ for any $t\in \mathbb{R}$.

Consider the one-parameter subgroups $\{\Theta(\operatorname{exp}(t\operatorname{ad}X)),t\in \mathbb{R}\}$ and $\{\operatorname{exp}(-t\operatorname{ad}X)$, $t\in \mathbb{R}\}$. 
Since $\Theta(y)=y^{-1}$, we have $\Theta(\operatorname{exp}(\operatorname{ad}X))=\operatorname{exp}(-\operatorname{ad}X)$.
That is, these two one-parameter subgroups coincides at $t=1$, and thus are equal.
By taking differential at $t=0$, we obtain that $\theta(X)=-X$ and $X\in \mathfrak{p}$.

Let $k=x\operatorname{exp}(-\frac{1}{2}\operatorname{ad}X)$.
Then 
\begin{align*}
    \Theta(k)^{-1}k&=\operatorname{exp}(-\frac{1}{2}\operatorname{ad}X) \Theta(x)^{-1}x \operatorname{exp}(-\frac{1}{2}\operatorname{ad}X)\\
    &=\operatorname{exp}(-\frac{1}{2}\operatorname{ad}X)\operatorname{exp}(\operatorname{ad}X)\operatorname{exp}(-\frac{1}{2}\operatorname{ad}X)=1.
\end{align*}
Thus $\Theta(k)=k$ and $k\in \mathbf{K}$.
So $x=k\operatorname{exp}(\frac{1}{2}\operatorname{ad}X)$ is the image of $(k,\frac{1}{2}X)\in \mathbf{K}\times \mathfrak{p}$ and the map is surjective.

Now we show that the map is injective.
By the previous argument, any $x\in \mathbf{G}$ can be written of the form $k\operatorname{exp}(\operatorname{ad}X)$ for some $k\in \mathbf{K}$ and $X\in \mathfrak{p}$. 
Repeating the same process as above, the eigenvectors of $\Theta(x)^{-1}x$ corresponding to the eigenvalues $\lambda_1,\cdots,\lambda_m$ form a basis of $\mathfrak{g}$.
Then with respect to this basis,
\begin{align*}
   X=\operatorname{diag}\{\operatorname{ln}\lambda_1,\cdots,\operatorname{ln}\lambda_m\}
\end{align*}
 is uniquely determined.
Thus $k$ is also uniquely determined and the map is injective.

Obviously the map is smooth. For the smoothness of its inverse map, we refer the proof of \cite[Theorem 6.31]{Knapp} for details.
The proof is finished.
\end{proof}

The diffeomorphism in the above proposition is called the Cartan decomposition of $\mathbf{G}$, and $\Theta$ is called the Cartan involution of $\mathbf{G}$.




\subsection{Iwasawa Decomposition}
In this subsection, we study the Iwasawa decomposition for the Lie algebra $\mathfrak{g}$ and the Chevalley group $\mathbf{G}$.

Let $\mathfrak{a}$ be the $\mathbb{R}$-span of $H_X,X\in \mathcal{R}$. 
Then $\mathfrak{a}$ is a maximal abelian subspace of $\mathfrak{p}$.

\begin{remark}
    For $\lambda\in\mathfrak{a}^*$, let 
    \begin{align*}
        \mathfrak{g}_{\lambda}=\{X\in \mathfrak{g}|[H,X]=\lambda(H)X, \forall H\in \mathfrak{a}  \}.
    \end{align*}
    If $\lambda\neq 0$ and $\mathfrak{g}_{\lambda}\neq \{0\}$, then $\lambda$ is called a restricted root of $\mathfrak{g}$.
    The set of restricted roots $\Sigma$ form an abstract root system in $\mathfrak{a}^*$.
    In general, $\Sigma$ is not equal to $\Phi$.
    However, in our case, the restricted roots corresponds bijectively to the roots of $\mathfrak{g}$, and the two root systems are equal.
\end{remark}

Now we arbitrarily fix a complete section of $\mathcal{R}$, and get a hereditary subcategory $\mathcal{B}$ of $\mathcal{R}$.
We fix a complete set of representatives $\{S_1,\cdots,S_n\}$ of isomorphism classes of simple objects in $\mathcal{B}$.

Let $\mathfrak{u}_+$ be the Lie subalgebra of $\mathfrak{g}$, $\mathbb{C}$-spanned by $u_X,X\in \operatorname{ind}\mathcal{B}$.
Then $\mathfrak{u}_+$ is the Lie algebra of $\mathbf{U}^+$.

\begin{proposition}
    We have an $\mathbb{R}$-vector space decomposition
    \begin{align*}
        \mathfrak{g}=\mathfrak{r}\oplus \mathfrak{a} \oplus \mathfrak{u}_+.
    \end{align*}
    This is called the Iwasawa decomposition for $\mathfrak{g}$.
\end{proposition}

\begin{proof}
    Notice that $\mathfrak{r}$ has an $\mathbb{R}$-basis
    \begin{align*}
         \{iH_{S_1},\cdots,iH_{S_n};u_X+u_{TX},i(u_X-u_{TX}),X\in \operatorname{ind}\mathcal{R}\},
    \end{align*}
    $\mathfrak{a}$ has an $\mathbb{R}$-basis $\{H_{S_1},\cdots,H_{S_n}\}$, and $\mathfrak{u}_+$ has an $\mathbb{R}$-basis $\{u_X,iu_X,X\in \operatorname{ind}\mathcal{B}\}$.
     The proposition follows.
\end{proof}

Let $\mathbf{A}$ be the subgroup of $\mathbf{G}$ generated by $h_{[X]}(t),X\in\operatorname{ind}\mathcal{R},t\in \mathbb{R}_{>0}$.
Then it is simply-connected and has Lie algebra $\mathfrak{a}$.


\begin{theorem}
    The multiplication $\mathbf{K}\times \mathbf{A}\times \mathbf{U}^+ \rightarrow \mathbf{G}, (k,a,u)\mapsto kau$ is a diffeomorphism.
    This is called the Iwasawa decomposition for $\mathbf{G}$.
\end{theorem}

\begin{proof}
    By definition, $\mathbf{A}$ is a closed subgroup of $\mathbf{H}$.
    Using the properties of $\mathbf{B}^+=\mathbf{U}^+\mathbf{H}$, the multiplication $\mathbf{A}\times \mathbf{U}^+\rightarrow \mathbf{A}\mathbf{U}^+$ is bijective.
    The subgroup $\mathbf{A}\mathbf{U}^+$ is closed, and has Lie algebra $\mathfrak{a}\oplus \mathfrak{u}_+$.
    Considering differentials, the multiplication $\mathbf{A}\times \mathbf{U}^+\rightarrow \mathbf{A}\mathbf{U}^+$ is a diffeomorphism.

    We can regard the multiplication map $\mathbf{K}\times \mathbf{A}\times \mathbf{U}^+ \rightarrow \mathbf{G}$ as the composition of maps $\mathbf{K}\times \mathbf{A}\times \mathbf{U}^+ \rightarrow \mathbf{K}\times \mathbf{A} \mathbf{U}^+ \rightarrow \mathbf{G} $.
    The image is the product of a compact subset and a closed subset, and is thus closed.
    On the other hand, by the Iwasawa decomposition of $\mathfrak{g}$, the differential $\mathfrak{r}\times (\mathfrak{a}\oplus \mathfrak{u}_+)\rightarrow \mathfrak{g},(X,Y)\mapsto X+Y$ is an isomorphism.
    Then the image of the multiplication map is open by the submersion theorem.
    Since $\mathbf{G}$ is connected, the image $\mathbf{K} \mathbf{A} \mathbf{U}^+=\mathbf{G}$ and the multiplication map is surjective.

    By definition, we have $\Theta(\mathbf{A} \mathbf{U}^+)=\mathbf{A} \mathbf{U}^-$, and $\Theta(h)=h^{-1}$, for any $h\in \mathbf{A}$.
    Since $\mathbf{A} \mathbf{U}^+ \bigcap \mathbf{A} \mathbf{U}^-=\mathbf{A}$, the only $\Theta$-fixed point in $\mathbf{A} \mathbf{U}^+$ is 1.
    We have shown that $\mathbf{K}=\mathbf{G}^{\Theta}$, and thus $\mathbf{K}\bigcap \mathbf{A} \mathbf{U}^+ =\{1\}$.
    This implies that the multiplication map is injective.
    Then it is easy to see that the multiplication map is a diffeomorphism and the theorem is proved.
\end{proof}


\begin{thebibliography}{99}
\bibitem[1997ROOT]{1997ROOT} L.~Peng and J.~Xiao. \newblock Root categories and simple lie algebras. \newblock {\em Journal of Algebra}, 198(1):19--56, 1997.
\bibitem[Knapp]{Knapp} A.~W. Knapp. \newblock {\em Lie groups beyond an introduction}, volume 140 of {\em Progress in Mathematics}. \newblock Birkh\"auser Boston, Inc., Boston, MA, second edition, 2002.
\bibitem[notes2]{notes2} B.~Li and J.~Xiao. \newblock Notes on {C}hevalley groups and root category {II}: Compact {L}ie groups and representations. \newblock {\em Journal of Algebra}, 711:113--155, 2027.
\end{thebibliography}
\end{document}
